\section{Introduction}

The canonical Arthur filtration organizes unramified automorphic functions
by nilpotent orbits in the Langlands dual Lie algebra. Its definition uses
singular support in a category of $\ell$-adic sheaves, whereas the underlying
space of finitely supported functions has a natural rational form. The
rationality question asks whether these two structures are compatible.

Arthur's conjectures organize departures from temperedness by an
additional algebraic $\mathrm{SL}_2$ factor in an automorphic parameter
\cite[Sections~4 and~8]{Arthur1989}. In geometric Langlands, Arinkin--Gaitsgory
introduced nilpotent singular support to encode this Arthur direction
\cite[Sections~1.1.2--1.1.6]{ArinkinGaitsgorySingularSupport}.
The restricted local-system stack and nilpotent automorphic category of
Arinkin--Gaitsgory--Kazhdan--Raskin--Rozenblyum--Varshavsky provide the
corresponding setting over a finite field \cite{AGKRRV1}.

Gaitsgory--Lafforgue--Raskin formulate this question, jointly with Kazhdan,
as the assertion that their canonical filtration is defined over $\Q$
\cite[Section~2.7, Conjecture~2.7.3]{GLR}. The construction and its relation
to Arthur's and Ramanujan's conjectures are discussed in
\cite{GLR,Raskin}. Restricted geometric Langlands theory supplies the
spectral description of the support categories, and the trace-of-Frobenius
theorem identifies the ambient trace with automorphic functions
\cite{AGKRRV1,AGKRRV3,GR}. These results make rationality a precise question
about the images of supported categorical traces.

We prove the rationality assertion for split connected semisimple groups
under the explicit characteristic assumptions below. Thus the result
positively resolves Conjecture~2.7.3 of \cite{GLR} in this setting. It
applies to the entire space of finitely supported functions, including its
noncuspidal part.

\subsection{Hypotheses and the filtration}

Let $X$ be a smooth projective geometrically connected curve over
$\mathbb F_q$, and let $G/\mathbb F_q$ be split, connected and semisimple.
Put $k_0=\overline{\mathbb F}_q$, $H=G_{k_0}$ and
$\mathfrak h=\Lie(H)$.

\begin{assumption}\label{intro:characteristic}
The following four conditions hold.
\begin{enumerate}[label=\textup{(\roman*)}]
\item The Lie algebra $\mathfrak h$ has an $H$-invariant nondegenerate
symmetric bilinear form whose restriction to the center of $\Lie(M)$ is
nondegenerate for every Levi subgroup $M$ of $H$.
\item For each Levi subgroup $M$ of $H$, including $H$ itself, and each
maximal torus $T_M\subset M$, restriction induces an isomorphism
\[
k_0[\mathfrak m]^M\xrightarrow{\ \sim\ }
k_0[\mathfrak t_M]^{W_M},\qquad
\mathfrak m=\Lie(M),\quad \mathfrak t_M=\Lie(T_M),\quad
W_M=N_M(T_M)/T_M.
\]
\item The scheme-theoretic centralizer in $H$ of every semisimple element
of $\mathfrak h$ is a Levi subgroup.
\item For every extension $K/k_0$ and every nilpotent
$n\in\mathfrak h\otimes_{k_0}K$, there is a parabolic subgroup
$P\subset H_K$, defined over $K$, such that
$n\in\Lie(R_u(P))$.
\end{enumerate}
Here a Levi subgroup is a Levi factor of a parabolic subgroup. Semisimple
means geometrically conjugate into the Lie algebra of a maximal torus;
nilpotent means that the geometric adjoint orbit closure contains zero.
\end{assumption}

Fix a prime $\ell\ne\operatorname{char}(\mathbb F_q)$ and set $E=\E$.
Write $S$ for the set of isomorphism classes of $G$-bundles on $X$ over
$\mathbb F_q$. For a characteristic-zero field $L$, let
\[
\cA_L=L^{(S)}
\]
be the space of finitely supported functions on $S$. Its distinguished
basis consists of the point functions. In particular
$\cA_E=E\otimes_\Q\cA_\Q$.

Fix the pinned split $\Q$-form of the Langlands dual group and write
$\Ghat/E$ for its scalar extension, with Lie algebra $\ghat$
and nilpotent cone $\Nhat$. All later coefficient identifications use
this same pinned form. On
$\overline X=X\times_{\mathbb F_q}k_0$ put
\[
\cC=\Shv_{\Nilp}(\Bun_G(\overline X),E),\qquad
\Phi=(\Frob_{\Bun})_*.
\]
The subscript denotes singular support in the global nilpotent cone, the
zero fiber of the Hitchin map. We use geometric Frobenius and its
pushforward on the automorphic category.

For a closed $\Ad(\Ghat)$-invariant subset $Y\subset\Nhat$, let
$\cC_Y\subset\cC$ be the derived-Satake support category. More explicitly,
the restricted geometric Langlands equivalence identifies
\begin{equation}\label{intro:restricted-equivalence}
\cC\simeq\IndCoh_{\Nhat}(Z),\qquad
Z={}^{\prime}\!\LocSys_{\Ghat}^{\restr}(\overline X),
\end{equation}
and $\cC_Y$ corresponds to $\IndCoh_Y(Z)$. The prime denotes the union of
connected components occurring in that equivalence; it is retained
throughout. A singular point is a pair $(\sigma,A)$ with $A$ a horizontal
section of the adjoint local system, after identifying $\ghat^*$ with
$\ghat$ by an invariant form. The support condition requires $A$ to take
values in $Y$. Thus $Y$ is a condition on singular directions, not on the
support of a function on $S$.

For a dualizable differential graded category $\cD$ and an endofunctor
$T$, write $\Tr(T,\cD)$ for its categorical trace, a complex of vector
spaces. It is Hochschild homology with coefficients in the graph bimodule
of $T$. The trace/functions isomorphism gives
$\Tr(\Phi,\cC)\simeq\cA_E$, with the latter in degree zero. Define
\begin{equation}\label{intro:filtration}
\cF_Y\cA_E=
\im\bigl[H^0\Tr(\Phi,\cC_Y)\longrightarrow\cA_E\bigr].
\end{equation}
This image definition does not assume concentration of the supported
trace. The required concentration and injectivity statements will be
proved in Section~\ref{tr:section}.

\begin{theorem}\label{intro:main}
Under Assumption~\ref{intro:characteristic}, for every prime
$\ell\ne\operatorname{char}(\mathbb F_q)$ and every closed invariant
$Y\subset\Nhat$, the natural map
\begin{equation}\label{intro:main-map}
E\otimes_\Q\bigl(\cA_\Q\cap\cF_Y\cA_E\bigr)
\xrightarrow{\ \sim\ }\cF_Y\cA_E
\end{equation}
is an isomorphism. The intersection is taken inside $\cA_E$.
\end{theorem}

No Hecke-finiteness or cuspidality condition is imposed. The theorem does
not require $X(\mathbb F_q)$ to be nonempty. These rational filtration
subspaces are also independent of $\ell$, with the closed invariant
supports matched through the pinned split dual group, as shown in
Section~\ref{desc:section}. On the unramified cuspidal subspace,
Corollary~\ref{desc:cuspidal-comparison} further identifies this filtration
with the sums of the rational cuspidal Arthur summands of
\cite[Theorem~1.1]{CuspidalArthur} whose orbit labels lie in $Y$.

\subsection{A bilinear detection criterion}

The proof gives a characterization after every abstract field
isomorphism $\iota:E\xrightarrow{\sim}\C$. Such isomorphisms are used
without prescribing their restriction to the algebraic numbers.
Transport all coefficient-linear categories and spaces by $\iota$.

Fix a closed point $v$ of $X$ of degree $d$, and put $r=q^d$. For a
nilpotent orbit $o\subset\ghat_\C$, choose an $\mathrm{SL}_2$-homomorphism
associated with a triple for $e\in o$, and write $\lambda_e$ for its
diagonal cocharacter, so $e$ has weight two. Define
\begin{equation}\label{intro:compactsets}
B_o=\bigl\{[\lambda_e(\sqrt r)s]:
s\in Z_{\Ghat_\C}(\mathrm{SL}_{2,e})_{\cpt}\bigr\}
\subset\Ghat_\C\sslash\Ghat_\C,
\qquad
B_Y=\bigcup_{o\subset Y_\C}B_o.
\end{equation}
Here $\sqrt r$ is positive, $[\ ]$ denotes semisimple conjugacy value,
and the compact form includes the components of the triple centralizer.
The sets $B_o$, as in \cite[Sections~3.1.1--3.1.3]{GLR}, do not depend
on these choices, are compact and pairwise disjoint. Their elementary properties are proved in
Lemma~\ref{desc:compactsets}.

For a finite-dimensional algebraic representation $D$ of $\Ghat_\C$,
let $T_D$ be the spherical Hecke operator at $v$, with the positive
square-root Satake normalization. For $f,g\in\cA_\C$, put
\begin{equation}\label{intro:pairing}
[f,g]=\sum_{b\in S}\frac{f(b)g(b)}{|\Aut(b)|}.
\end{equation}
This is a bilinear pairing. The automorphism groups over the finite field
are finite, and the sum is finite.

The key result, Theorem~\ref{meas:detector}, says that for each $f,g$
there is a unique finite complex Borel measure $\nu_{f,g}$ on
$B=\bigcup_oB_o$ such that
\begin{equation}\label{intro:moments}
[T_Df,g]=\int_B\chi_D\,d\nu_{f,g}
\quad\text{for every }D,
\end{equation}
and that
\begin{equation}\label{intro:detector}
f\in\iota_*(\cF_Y\cA_E)
\quad\Longleftrightarrow\quad
\supp(\nu_{f,g})\subseteq B_Y
\quad\text{for every }g\in\cA_\C.
\end{equation}
A finite complex measure means one of finite total variation. No positivity
of these measures is asserted or needed.

The right side of \eqref{intro:detector} uses the same complex function
space, point-counting Hecke operations, rational stack weights and compact
sets for every $\iota$. It therefore gives invariance of the filtration
under all coefficient automorphisms over $\Q$. Finite-coordinate linear
algebra then proves \eqref{intro:main-map}.

\subsection{The local arguments}

The supported trace concentration and injection statements in
\cite[Corollaries~2.1.10 and~2.1.13]{GLR} are conditional on
Conjecture~2.1.5 there, as stipulated in its Section~2.1.9.
Section~\ref{tr:section} proves the concentration and injection needed
here from the component and orbit calculations, and constructs actual
compact Weil representatives of the supported trace classes.

The central task is to detect every nonzero top support class by a
finitely supported function test despite possible cancellation from its
boundary.
Our contribution is the bilinear criterion
\eqref{intro:moments}--\eqref{intro:detector} on the full space of
finitely supported functions, together with the local calculations and
compact representatives that establish it. First, on a retained fixed
component of the local-system stack, a pure Weil lift supplies separated
Frobenius weights. Finite covariant tests recover the component from a
closed-orbit slice while preserving completion along its invariant base.
Nonresonance then gives a quadratic moment-map model compatible with
evaluation bundles and singular support.

A deformation-theoretic antecedent is the quadratic-germ theorem of
Goldman--Millson for compact K\"ahler manifolds and compact-image
representations \cite[Theorem~1]{GoldmanMillson}. Pridham obtains
Frobenius-compatible classical deformation hulls and cup-product
quadratic equations from purity for smooth proper varieties
\cite[Theorem~2.11 and Corollary~2.12]{Pridham}. Our finite-jet
nonresonance argument retains the derived equivariant model, its formal
invariant base, evaluation bundles and singular-support map.

Second, a graded Koszul description reduces each nilpotent-orbit
quotient to equivariant Clifford modules together with strictly
contracting Hom directions. Clifford algebras and graded matrix
factorizations also enter the spectral Whittaker construction of
Ben-Zvi, Raskin and Venkatesh described in \cite[Section~4.5]{Raskin},
which develops an idea of V.~Lafforgue; the even-nilpotent construction
appears in \cite[Section~18.5]{BZSV}. Here a twisted-bar contraction
proves trace concentration. Proper orbit-closure extensions and finite Koszul lifts
produce actual compact Weil classes, so the result describes the image
in \eqref{intro:filtration}, not only an abstract associated graded.
The passage from equivariant localization to a cofiber sequence of
traces follows the formalism of Hoyois--Scherotzke--Sibilla
\cite[Theorem~3.4 and Proposition~5.4]{HSS}.

Third, nonstandard automorphic duality turns these classes into the
bilinear tests \eqref{intro:pairing}. Localization with supports gives
decaying normal-direction contributions and hence finite complex
measures. On a top orbit the common density is invertible; finite
character approximations remove it and leave a nondegenerate coefficient
pairing. A strict drop in ambient adjoint rank separates every proper
boundary contribution. These two facts prevent cancellation from hiding
a nonzero top class.

The supported local-cohomology and orbit-distribution strategy is inspired
by Kazhdan--Okounkov's treatment of the unramified Eisenstein spectrum
\cite[Sections~2.6.1--2.6.4 and~3.3--3.4]{KazhdanOkounkov}.
The all-function bilinear detector, compact Weil lifts, and
noncancellation and rational-descent arguments are developed here;
that reference does not establish the rationality theorem above.

\begin{figure}[ht]
\centering
\begin{tikzcd}[row sep=2em]
\text{quadratic component model}\quad
\text{\footnotesize(Section~\ref{mod:section})}
  \arrow[d,"\text{orbit traces}"] \\
\text{compact Weil trace classes}\quad
\text{\footnotesize(Sections~\ref{orb:section}--\ref{tr:section})}
  \arrow[d,"\text{supported Hom pairings}"] \\
\text{bilinear measure criterion}\quad
\text{\footnotesize(Section~\ref{meas:section})}
  \arrow[d,"\text{coefficient descent}"] \\
\text{rational form of }\cF_Y\cA_E\quad
\text{\footnotesize(Section~\ref{desc:section})}
\end{tikzcd}
\caption{Main line of the proof. The compact-set properties used by
the measure criterion are proved independently in
Section~\ref{desc:section}.}
\label{intro:diagram}
\end{figure}

Section~\ref{inp:section} records the established geometric inputs and
their precise scope. Sections~\ref{mod:section}--\ref{tr:section} prove the
local model and trace calculations. Section~\ref{meas:section} establishes
the measure criterion, and Section~\ref{desc:section} completes the
rational descent.
